\section{Introduction}
Perfect correlation with human mean ratings does not identify a judge's rating distribution.
A conditional distribution describes how human ratings vary for an item.
Its mean summarizes only one aspect of that variation.
If two forecasts have identical means on every item, correlation of those means assigns identical values to both forecasts.
The distributions can still disagree about uncertainty and the probabilities of individual ratings.

The interpretation of correlation therefore depends on the prediction target.
When the target is association between scalar scores, correlation directly describes that association.
When the target is a conditional rating distribution, association alone leaves additional questions unanswered.
Strictly proper scoring rules address distribution recovery by making the true distribution the unique optimum of expected score~\cite{gneiting2007}.
The ranked probability score uses the ordering of the rating categories~\cite{epstein1969}.
Our comparison concerns these mathematical targets rather than a claim about how often practitioners use any objective.

We characterize the information discarded by Pearson correlation of conditional means.
At correlation one, forecast means equal target means after a common positive affine transformation.
On a scale with at least three categories, unrestricted distributions also retain shape variation at fixed means.
This result concerns the objective's dependence on outputs.
It places no linearity restriction on a neural judge's map from an item to its mean rating.

The model class determines which ambiguities are feasible.
A restricted family can tie shape to the mean or exclude every alternative affine copy of the target means.
Shared parameters can also couple changes in mean and shape.
We account for these restrictions through the model's mean map and its feasible directions.
We then compare population cross-entropy, normalized individual-rating ranked probability score, and total variation distance to the true distribution.

We give detailed proofs and an exact finite illustration rather than a trained-model comparison.
The illustration holds the means fixed across three strictly positive rating distributions.
Its correlation is one, although its average normalized excess ranked probability score is $3/200$.
The supplied arithmetic reproduces these values without fitting a model.
A separate analysis explains why finite rating noise and estimation error prevent a universal training comparison.

The results support matching the evaluation criterion to the intended output.
They do not establish that one training algorithm always generalizes better.
We make the following contributions:
\begin{enumerate}
\item We characterize the invariances, equality conditions, gradients, and probability-profile variants of Pearson correlation.
\item We derive distributional loss identities, moment bounds, and lower bounds for recovery from means alone.
\item We state how the rating scale, model restrictions, and finite-data assumptions change these conclusions.
\end{enumerate}
